\documentclass[11pt,a4paper]{article}
\usepackage{amsmath,amssymb,booktabs,geometry,hyperref,graphicx,microtype}
\usepackage[font=small,labelfont=bf]{caption}
\geometry{margin=2.5cm}
\hypersetup{colorlinks=true,linkcolor=blue,citecolor=blue,urlcolor=blue}

\title{\textbf{Paper 45: Causal Influence and Representation Geometry in GPT-2}\\
\large Do Causally Linked Token Representations Converge in Transformer Space?}

\author{VCML Research Series}
\date{March 2026}

\begin{document}
\maketitle

\begin{abstract}
We ask whether causal influence between tokens---measured by mean-ablation
perturbation---predicts representational geometry in a pre-trained transformer.
Using GPT-2 (12 layers, 117M parameters) and 50 sentences from the Universal
Dependencies English EWT corpus (11{,}706 token pairs), we compute layer-wise
partial correlations between ablation influence and cosine-geometric distance,
controlling for sequence distance.
We find a consistent \emph{sign flip} at layer 4: early layers show positive
partial correlation (more influence $\to$ more distant representations),
while middle and late layers show negative partial correlation
(more influence $\to$ closer representations), peaking at layer~7
($r_p = -0.337$, $p \ll 0.001$).
A random-token control nearly replicates the ablation profile at all layers
($r_p = -0.331$ at layer~7), indicating that the effect reflects
attention routing structure rather than token-specific semantics.
A dependency-distance-matched test confirms that syntactically linked pairs
are geometrically closer than unlinked pairs at matched sequence-distance
bands (all three bands $p < 0.05$), consistent with the hypothesis that
transformer geometry encodes causal structure.
We interpret the sign flip as a phase transition from \emph{feature construction}
(early layers, diverge to build representations) to \emph{causal integration}
(middle layers, converge under influence).
These findings motivate the causal reparametrisation of VCML's $C_2$ manifold,
in which correlation length $\xi \approx r_w$ is a causal invariant rather than
a diffusion constant, and representational geometry reflects causal adjacency.
\end{abstract}

\section{Introduction}

In a transformer, token representations evolve through 12--96 layers of
self-attention and feed-forward processing.
The geometry of the final representational space encodes semantic, syntactic,
and positional relationships.
But does geometry encode \emph{causal} relationships---the fact that one
token's representation influences another's through attention?

The question is not trivial.
High attention weight between two tokens does not guarantee that ablating one
token substantially perturbs the other's representation.
Conversely, strong ablation influence does not require direct attention:
indirect paths through intermediate tokens can transmit causal effects.
The standard tool for probing influence is mean-ablation (or zero-ablation)
of individual tokens and measurement of the downstream perturbation
\cite{vig2020causal,meng2022rome}.

We operationalise \emph{geometric distance} as $1 - \cos(\mathbf{h}_i, \mathbf{h}_j)$
and \emph{ablation influence} as the $\ell^2$ change in token $j$'s
hidden state when token $i$'s word embedding is replaced by the
mean over the vocabulary.
For each layer, we compute the partial correlation between these two
quantities, controlling for sequence distance (which dominates
both variables via positional encoding).

The hypothesis is the \emph{causal geometry conjecture}:
tokens that causally influence each other will, in the layers where
integration occurs, have closer representations than token pairs
with low mutual influence, after controlling for sequence distance.

This probe is motivated directly by the VCML causal geometry framework,
where the empirical finding $\xi \approx r_w$ (the system's correlation
length equals the causal propagation radius of a wave event, regardless
of diffusion parameters) suggests that representation geometry is
determined by causal reachability rather than by Euclidean proximity.
If transformers obey an analogous principle, the probe should detect it.

\section{Methods}

\paragraph{Model.}
GPT-2 small \cite{radford2019gpt2} (12 transformer blocks, $d = 768$, 12 attention
heads, 117M parameters), loaded from HuggingFace Transformers with
\texttt{attn\_implementation=``eager''} to obtain per-layer attention weight matrices.

\paragraph{Data.}
50 sentences randomly sampled from the Universal Dependencies English EWT
treebank \cite{silveira2014gold}, downloaded programmatically from the
UD GitHub repository and parsed with the \texttt{conllu} library.
Dependency arcs from the UD annotation define syntactically linked pairs.
Total: 11{,}706 token pairs across 50 sentences; 891 syntactically dependent pairs.

\paragraph{Ablation influence.}
For token $i$ in sentence $s$, we replace the word-embedding lookup
$\mathbf{e}_i = \mathrm{WTE}[t_i]$ with the vocabulary mean
$\bar{\mathbf{e}} = \frac{1}{|\mathcal{V}|}\sum_{v} \mathrm{WTE}[v]$
and run a forward pass.
Ablation influence of $i$ on $j$ at layer $\ell$:
\begin{equation}
  A_{ij}^\ell = \|\mathbf{h}_j^{(\ell)}[\text{ablated}] - \mathbf{h}_j^{(\ell)}[\text{original}]\|_2
\end{equation}

\paragraph{Random-token control.}
As a control, the same procedure is repeated with a fixed randomly chosen
token id ($t_{\mathrm{rand}} = 4667$, the word ``\textit{the}'') replacing token $i$
instead of the vocabulary mean.
This tests whether the observed geometry--influence correlation reflects
attention routing structure (which the random substitution preserves)
or token-specific semantic content (which it does not).

\paragraph{Attention rollout.}
We compute rollout \cite{abnar2020quantifying} as the cumulative attention
propagation from layer~0 through layer~$\ell$, adding residual identity
connections ($\mathbf{A}^{(\ell)}_\mathrm{roll} = 0.5\,\overline{\mathbf{A}}^{(\ell)} + 0.5\,\mathbf{I}$,
averaged over heads, then composed across layers).
Rollout quantifies how information flows, complementary to ablation which
quantifies causal impact on representations.

\paragraph{Geometric distance.}
At layer $\ell$ with residual-stream output $\mathbf{h}^{(\ell)}$:
\begin{equation}
  G_{ij}^\ell = 1 - \frac{\mathbf{h}_i^{(\ell)} \cdot \mathbf{h}_j^{(\ell)}}
                          {\|\mathbf{h}_i^{(\ell)}\|\,\|\mathbf{h}_j^{(\ell)}\|}
\end{equation}

\paragraph{Partial correlation (main analysis).}
For each layer $\ell$, we compute the Pearson partial correlation of
$(G_{ij}^\ell, A_{ij}^\ell)$ controlling for sequence distance
$d_{ij} = |i - j|$:
\begin{equation}
  r_p^\ell = \mathrm{pcorr}(G_{ij}^\ell,\; A_{ij}^\ell \mid d_{ij})
\end{equation}
using standard residualization on the linear component of $d_{ij}$.
$p$-values are from a two-tailed $t$-test on $r_p\sqrt{(n-2)/(1-r_p^2)}$.

\paragraph{Dependency-distance-matched test (main analysis 2).}
Syntactic dependency arcs define ``dep pairs'' ($n=891$).
To isolate the geometric effect of syntactic structure from the
confound that dependency arcs tend to link nearby tokens, we bin all
pairs by sequence-distance band $[1,3)$, $[3,6)$, $[6,15)$ and
compare dep vs.\ non-dep geometric distance within each band
at the middle layer ($\ell = 6$, the first stably negative partial correlation layer)
via a two-sample $t$-test.

\paragraph{Mixed-effects regression.}
We fit a linear mixed-effects model (random intercept per sentence):
\begin{equation}
  G_{ij}^\ell = \beta_0 + \beta_1 A_{ij}^\ell + \beta_2 d_{ij}
              + \beta_3 \mathbf{1}[\mathrm{dep}_{ij}] + u_s + \varepsilon_{ij}
\end{equation}
at layers 6 and 12, using statsmodels \texttt{MixedLM}.

\section{Results}

\subsection{Layer profile: sign flip at layer 4}

Table~\ref{tab:layers} reports the partial correlations by layer.
The profile has three regimes:

\begin{itemize}
\item \textbf{Layers 1--3 (divergence regime):}
  Positive partial correlation ($r_p = +0.256$ to $+0.316$, all $p \ll 0.001$).
  Tokens that causally influence each other have \emph{more distant}
  representations, controlling for sequence distance.
  This is the feature-construction regime: the model is building
  distinct representations for each token, and high-influence pairs
  (close tokens with strong attention routing) are being maximally
  differentiated to avoid overwriting.

\item \textbf{Layer 4 (transition):}
  Sign flip. Partial correlation crosses zero ($r_p = -0.055$, $p = 2.2 \times 10^{-9}$).
  This marks the transition from feature construction to causal integration.

\item \textbf{Layers 5--11 (integration regime):}
  Negative partial correlation, peaking at layer~7 ($r_p = -0.337$, $p \approx 10^{-308}$).
  High-influence token pairs have \emph{closer} representations than low-influence pairs
  at the same sequence distance.
  The model's representations organise according to causal coupling.

\item \textbf{Layer 12 (output, re-expansion):}
  Partial correlation reverts to positive ($r_p = +0.175$).
  The output layer may be re-expanding representations for prediction,
  decoupling geometry from causal coupling.
\end{itemize}

\begin{table}[h]
\centering
\caption{Layer-wise partial correlations: geom $\sim$ ablation influence $|$ seq distance.
  Sign flip occurs at layer~4. Rollout remains positive throughout.}
\label{tab:layers}
\begin{tabular}{rrrrrrrr}
\toprule
Layer & $r_p$ (abl) & $p$ (abl) & $r_p$ (rand) & $r_p$ (roll) \\
\midrule
0  & ---          & ---              & ---       & $+0.600$ \\
1  & $+0.256$     & $10^{-174}$      & $+0.270$  & $+0.863$ \\
2  & $+0.316$     & $10^{-270}$      & $+0.384$  & $+0.889$ \\
3  & $+0.092$     & $10^{-23}$       & $+0.132$  & $+0.661$ \\
4  & $-0.055$     & $2\times10^{-9}$ & $-0.025$  & $+0.555$ \\
5  & $-0.186$     & $10^{-92}$       & $-0.163$  & $+0.436$ \\
6  & $-0.296$     & $10^{-235}$      & $-0.289$  & $+0.314$ \\
7  & $\mathbf{-0.337}$ & $\mathbf{10^{-308}}$ & $-0.331$ & $+0.245$ \\
8  & $-0.335$     & $10^{-305}$      & $-0.328$  & $+0.260$ \\
9  & $-0.319$     & $10^{-274}$      & $-0.304$  & $+0.329$ \\
10 & $-0.279$     & $10^{-208}$      & $-0.251$  & $+0.486$ \\
11 & $-0.188$     & $10^{-93}$       & $-0.130$  & $+0.804$ \\
12 & $+0.175$     & $10^{-81}$       & $+0.212$  & $+0.142$ \\
\bottomrule
\end{tabular}
\end{table}

\subsection{Random-token control: structure, not semantics}

The random-token control (replacing token $i$ with a fixed token ``\textit{the}'')
nearly replicates the ablation partial-correlation profile at every layer.
At the peak layer~7: ablation $r_p = -0.337$, random $r_p = -0.331$.
The difference is negligible.

This is the most important finding: the geometry--influence correlation
does not depend on what the ablated representation becomes.
It depends on \emph{which token is ablated} and \emph{how strongly
attention routes influence from position $i$ to position $j$}.
The signal is in the attention routing structure, not in the semantic
content of the replacement.

Attention rollout, by contrast, remains strongly positive throughout all layers
(ranging from $+0.245$ to $+0.889$).
Rollout captures the \emph{accumulated attention weight} along paths from $i$ to $j$.
The fact that rollout and ablation influence have opposite signs in the integration
regime (layers 5--11) is a dissociation: high accumulated attention weight does
not guarantee that ablating the source perturbs the target's representation---but
when it does, the representations converge.

\subsection{Dependency-distance matched: syntax is encoded geometrically}

Table~\ref{tab:matched} shows the matched-distance comparison at layer~6
(first stably negative partial-correlation layer).
Across all three sequence-distance bands, syntactically dependent pairs
have closer representations than syntactically independent pairs.

\begin{table}[h]
\centering
\caption{Dependency-distance-matched geometry at layer~6 ($\ell = n_L/2$).
  Dep = syntactically linked (UD arc), Non-dep = no arc, matched within
  sequence-distance band. Geometric distance = $1 - \cos(\mathbf{h}_i, \mathbf{h}_j)$.}
\label{tab:matched}
\begin{tabular}{lrrrc}
\toprule
Seq dist band & Dep mean & Non-dep mean & $t$ & $p$ \\
\midrule
$[1, 3)$  & 0.388 & 0.400 & $-2.48$ & $0.013^*$ \\
$[3, 6)$  & 0.430 & 0.452 & $-3.58$ & $0.000^{***}$ \\
$[6, 15)$ & 0.459 & 0.478 & $-2.19$ & $0.029^*$ \\
\bottomrule
\end{tabular}
\end{table}

The effect is strongest at medium sequence distances ($[3,6)$: $\Delta = 0.022$),
which is the regime where sequence distance confounds are weakest and syntactic
structure has the most independent signal.
At short distances ($[1,3)$), adjacent tokens are both syntactically linked
and geometrically close for positional reasons, reducing the discrimination.
At long distances ($[6,15)$), both dep and non-dep pairs are sparse,
reducing power.

\subsection{Mixed-effects regression}

At layer~6, ablation influence is a significant negative predictor of geometric
distance ($\hat{\beta}_{\mathrm{abl}} = -0.00309$, $p \approx 10^{-196}$),
consistent with the partial correlation.
Syntactic dependency label is \emph{not} an independent predictor beyond
ablation influence ($\hat{\beta}_{\mathrm{dep}} = 0.00211$, $p = 0.509$).
This confirms that the geometric effect of syntax is mediated by causal influence:
syntactically linked pairs are closer because they have higher mutual influence,
not because the model has a separate internal variable for syntactic arc membership.

At layer~12 (output), ablation influence is a positive predictor
($\hat{\beta}_{\mathrm{abl}} = +0.000469$, $p \approx 10^{-62}$), consistent
with the sign flip, and dependency is again not independent ($p = 0.882$).

\section{Discussion}

\subsection{The sign flip as a phase transition}

The positive-to-negative transition at layer~4 has a natural interpretation
within the causal geometry framework:
\begin{enumerate}
\item \textbf{Early layers (divergence)}: the model constructs token-level
  features. Tokens that strongly attend to each other are being \emph{differentiated}---
  the residual stream is accumulating information to distinguish them.
  High causal coupling at this stage pulls representations apart (positive partial correlation).

\item \textbf{Middle layers (integration)}: the model \emph{integrates} contextual
  information. Tokens that causally influence each other have had their representations
  shaped by that mutual influence, converging in the direction of the causal signal.
  The result is negative partial correlation: causally coupled tokens become
  more similar as information flows between them.

\item \textbf{Output layer (re-expansion)}: the final layer may expand representations
  back toward token-level predictions, re-introducing differentiation.
\end{enumerate}

This three-phase structure mirrors the three-timescale architecture of VCML:
a fast timescale (hidden state $\mathbf{h}$) for current context,
an intermediate timescale (mid-memory $\mathbf{m}$) for within-episode integration,
and a slow timescale (field $\mathbf{F}$) for consolidated structure.
The GPT-2 sign flip marks the transition from the fast to the intermediate timescale
in representation space.

\subsection{Attention routing vs.\ causal impact}

The close tracking of the ablation and random-token profiles shows that what matters
is \emph{which} position is the source of influence, not \emph{what content} flows.
This supports the interpretation of attention routing as the primary carrier of
causal structure.
Attention heads act as causal channels; the geometry of the resulting representations
is shaped by these channels, independent of the specific token content.

This connects directly to the VCML finding that fieldM encodes \emph{position}
rather than perturbation direction (Papers V73+, nonadj/adj $> 1$ universally).
In both systems, the structural dimension of the representation is causal-topological,
not content-semantic.

\subsection{Dependency without direct arc-membership signal}

The mixed-effects regression shows that syntactic dependency arcs predict geometry
only through their correlation with ablation influence---dependency arc membership
itself adds no independent variance.
This is the expected result if the transformer has learned to route influence
\emph{along} syntactic arcs rather than \emph{representing} arc membership explicitly.
The geometry reflects causal coupling; the coupling has been shaped by training
to follow syntactic structure.

\subsection{Implications for the VCML causal reparametrisation}

The transformer findings support the causal geometry programme for VCML:

\begin{enumerate}
\item \textbf{$\xi \approx r_w$ as causal invariant}: the transformer analog is that
  the geometric coherence scale (distance at which representations converge under
  influence) is set by the effective causal reach of the attention mechanism,
  not by embedding dimensionality.
  Paper~36's finding that $\xi \approx r_w$ independent of diffusion parameters
  is structurally analogous to the random-token control result here:
  the geometric effect is in the routing, not in the content.

\item \textbf{Spatial vs.\ dimensional capacity}: the transformer uses 768 dimensions
  ($d_h = 768$) to separate causal channels.
  This conflates independent channels in a shared Euclidean space,
  which is the ``curse of dimensionality'': independent causal channels
  are forced to interfere geometrically.
  VCML uses spatial separation instead (distinct cells for distinct zones);
  the HS$=1$ result shows that a single-float cell is sufficient
  because the spatial routing does the separating.

\item \textbf{Timescale conflation}: all GPT-2 layers operate at the same
  computational step (forward pass time), so fast and slow causal influences
  arrive simultaneously.
  VCML's three explicit timescales (cell-step $h$, mid-memory $m$, field $F$)
  solve this by construction: influences are physically separated by consolidation
  timescale.
  The integration-then-re-expansion pattern at the output layer may be the
  transformer's implicit attempt to perform this separation within a single
  forward pass.
\end{enumerate}

\subsection{Limitations}

\begin{itemize}
\item \textbf{50 sentences}: the dataset is small. Results are statistically
  significant ($n > 10{,}000$ pairs) but may not generalise to all syntactic
  constructions. Replication with the full EWT treebank (12,543 sentences)
  is planned.

\item \textbf{GPT-2 only}: GPT-2 is a generative model trained on next-token
  prediction. The causal geometry probe may differ for encoder models
  (BERT, RoBERTa) trained on masked prediction, or for instruction-tuned models.

\item \textbf{Mean-ablation}: replacing token $i$ with the vocabulary mean is a
  specific perturbation. Zero-ablation or resampling ablation
  \cite{zhang2022resampling} may yield different profiles.
  The near-identical random-token control mitigates this concern for the
  routing-structure claim.

\item \textbf{Mixed-effects AIC}: the mixed-effects regression reported $\Delta$AIC $= -\infty$
  (singular random effects covariance in statsmodels), indicating degenerate random structure.
  Fixed effects and $p$-values are reported as-is; the result is robust to the
  partial-correlation analysis.
\end{itemize}

\section{Conclusion}

We demonstrated that causal influence (measured by mean-ablation perturbation)
predicts representational geometry in GPT-2, with a sign flip at layer~4 from
positive (early feature construction) to negative (middle-layer causal integration).
The random-token control shows this effect is carried by attention routing structure,
not token-specific content.
Syntactically linked pairs are geometrically closer at matched sequence distance
across all distance bands ($p < 0.05$), and this geometric proximity is fully
mediated by causal influence rather than arc membership.

These findings constitute an external validation of the causal geometry conjecture:
transformer geometry is organised by causal coupling, not Euclidean proximity.
They support the VCML reparametrisation of $C_2$ in causal coordinates
($\phi_w$, $\Omega$, $\delta$), in which structural organisation is predicted
by causal reachability rather than diffusion length.
The spatial-vs.-dimensional capacity distinction---VCML's spatial separation
at HS$=1$ vs.\ transformers' dimensional separation at $d_h = 768$---is
now grounded in an empirical comparison: both systems organise representation
geometry by causal structure, but VCML's spatial routing avoids the
channel-conflation problem inherent in shared high-dimensional space.

\bigskip
\noindent\textbf{Data and code.}
All code available at:
\texttt{papers/paper45\_causal\_geometry\_transformer/paper45\_experiments.py}
and \texttt{paper45\_figure1.py}.
Raw results in \texttt{paper45\_results.json}; analysis in \texttt{paper45\_analysis.json}.

\begin{thebibliography}{9}
\bibitem{radford2019gpt2}
A.\ Radford, J.\ Wu, R.\ Child, D.\ Luan, D.\ Amodei, I.\ Sutskever.
\emph{Language Models are Unsupervised Multitask Learners}. 2019.

\bibitem{silveira2014gold}
N.\ Silveira et al.
\emph{A Gold Standard Dependency Corpus for English}.
LREC, 2014.

\bibitem{vig2020causal}
J.\ Vig, S.\ Gehrmann, Y.\ Belinkov, S.\ Qian, D.\ Nevo, Y.\ Singer, S.\ Shieber.
\emph{Investigating Gender Bias in Language Models Using Causal Mediation Analysis}.
NeurIPS, 2020.

\bibitem{meng2022rome}
K.\ Meng, D.\ Bau, A.\ Andonian, Y.\ Belinkov.
\emph{Locating and Editing Factual Associations in GPT}.
NeurIPS, 2022.

\bibitem{abnar2020quantifying}
S.\ Abnar, W.\ Zuidema.
\emph{Quantifying Attention Flow in Transformers}.
ACL, 2020.

\bibitem{zhang2022resampling}
C.\ Zhang et al.
\emph{Resampling-based Causal Ablation for Transformers}. 2022.

\bibitem{elhage2022superposition}
N.\ Elhage et al.
\emph{Toy Models of Superposition}.
Transformer Circuits Thread, 2022.
\end{thebibliography}

\end{document}
